\section{Conclusion}
\label{sec:conclusion}

We have presented \Fulcrum, a topology-aware differential-privacy
allocation for federated learning, paired with \TADI, a four-channel
distributional inference attack that targets the gap left open by
content-only privacy mechanisms. The defense rests on a per-client
mutual-information bound that decomposes additively into a controllable
mechanism term — shrinkable by per-client noise allocation — and an
uncontrollable prior-coupling floor that DP-SGD cannot reduce. Under a
fixed utility budget, the closed-form balanced min-max optimal
allocation $\sigma_i^{*\,2} = a/(K^\star - \ell_i^\circ)$ strictly
dominates uniform DP-SGD whenever the structural-leverage scores are
non-uniform across clients, and degenerates gracefully to uniform
allocation in the symmetric case. Three principled leverage proxies —
SBM-derived group size, degree-based, and FedAvg-influence-derived
dataset size — cover the canonical hierarchical, decentralised, and
cross-silo deployment regimes.

\PLACE{Empirical headline (pending Stage~6 completion):
report the area-between-curves advantage on each setting, the
$\eta$-sweep gap pattern across the four topologies, and confirm the
correlation between predicted privacy bound $K^\star$ and measured
\TADI attack lift.}

The framework provides a deployable defense that is consistent with its
threat model, comes with a provable bound, and admits a measurable
privacy--utility trade-off curve. Open directions include amplified-RDP
tightening of Theorem~\ref{thm:per-client} under Poisson subsampling,
extension to active adversaries, and analytical leverage bounds for
the heuristic proxies.

\paragraph{Reproducibility.}
The implementation, configurations, sweep specifications, and analysis
pipeline are released at
\url{https://github.com/murtazahr/Fulcrum} under an open licence. Each
experiment is fully specified by a YAML configuration; the SQLite-backed
runner is idempotent under config-hash, supporting reproducible
re-execution and partial-result resumption.
